% META: {"id": "1SPE-TRIGO-EV-B", "chapitre": "1SPE-TRIGONOMETRIE", "type_objet": "evaluation", "version": "B", "duree_min": 55, "bareme_total": 20, "capacites_codes": ["C1","C2","C3","C4","C5"], "competences": ["calculer","raisonner","representer","communiquer"], "mode_creation": "ex_nihilo", "sources_inspiration": [], "parametres_sympy": {}, "fichier_tex": "chapitres/1SPE-TRIGONOMETRIE/evaluations/1SPE-TRIGO-EV-B.tex", "corrige_tex": "chapitres/1SPE-TRIGONOMETRIE/evaluations/1SPE-TRIGO-EV-B-corrige.tex", "status": "generated"}
% BEGIN-VERIFY
% from sympy import *
% x = symbols('x', real=True)
% assert 315 * pi / 180 == 7*pi/4
% assert 5*pi/3 * 180/pi == 300
% assert Rational(-13,6) + 2 == Rational(-1,6)
% assert cos(3*pi/4) == -sqrt(2)/2
% assert sin(4*pi/3) == -sqrt(3)/2
% assert simplify(sin(7*pi/12) - (sqrt(6)+sqrt(2))/4) == 0
% g = 3*sin(x) - 1
% assert g.subs(x, pi/2) == 2
% assert g.subs(x, 3*pi/2) == -4
% END-VERIFY

%---------------------------------------------------------------
% EVALUATION B (55 min) — Trigonometrie
% Bareme total : 20 points
%---------------------------------------------------------------

\section*{Evaluation — Trigonometrie (Version B)}
\textit{Duree : 55 min. Calculatrice interdite. Toutes les reponses doivent etre justifiees.}

\bigskip

%===============================================================
\subsection*{Exercice 1 \hfill (4 points) — C1}
%===============================================================

\begin{enumerate}
  \item Convertir $315°$ en radians. \hfill (1 pt)
  \item Convertir $\dfrac{5\pi}{3}$ en degres. \hfill (1 pt)
  \item Determiner la mesure principale de $-\dfrac{13\pi}{6}$. \hfill (1 pt)
  \item Placer sur le cercle trigonometrique le point associe a $\dfrac{3\pi}{4}$ et donner ses coordonnees. \hfill (1 pt)
\end{enumerate}

%===============================================================
\subsection*{Exercice 2 \hfill (4 points) — C2}
%===============================================================

\begin{enumerate}
  \item Donner les valeurs exactes de $\cos\!\left(\dfrac{4\pi}{3}\right)$ et $\sin\!\left(\dfrac{4\pi}{3}\right)$. \hfill (1,5 pt)
  \item Donner la valeur exacte de $\cos\!\left(\dfrac{11\pi}{6}\right)$. \hfill (1 pt)
  \item Sachant que $\sin x = \dfrac{5}{13}$ et $x \in \left]0;\dfrac{\pi}{2}\right[$, calculer $\cos x$. \hfill (1,5 pt)
\end{enumerate}

%===============================================================
\subsection*{Exercice 3 \hfill (4 points) — C3}
%===============================================================

\begin{enumerate}
  \item Calculer $\sin\!\left(\dfrac{7\pi}{12}\right)$ en ecrivant $\dfrac{7\pi}{12} = \dfrac{\pi}{3} + \dfrac{\pi}{4}$. \hfill (2 pts)
  \item Montrer que $\cos(2x) = 2\cos^2 x - 1$. \hfill (1 pt)
  \item Calculer $\cos^2\!\left(\dfrac{\pi}{12}\right)$. \hfill (1 pt)
\end{enumerate}

%===============================================================
\subsection*{Exercice 4 \hfill (4 points) — C4}
%===============================================================

\begin{enumerate}
  \item Resoudre dans $\mathbb{R}$ : $\sin x = \dfrac{\sqrt{3}}{2}$. \hfill (1,5 pt)
  \item Resoudre sur $[0;2\pi]$ : $\cos x = -\dfrac{\sqrt{2}}{2}$. \hfill (1,5 pt)
  \item Resoudre sur $[0;2\pi]$ : $\sin x \geqslant \dfrac{1}{2}$. \hfill (1 pt)
\end{enumerate}

%===============================================================
\subsection*{Exercice 5 \hfill (4 points) — C5}
%===============================================================

On pose $g(x) = 3\sin x - 1$.

\begin{enumerate}
  \item Quelle est la periode de $g$ ? La fonction est-elle impaire ? \hfill (1 pt)
  \item Dresser le tableau de variations de $g$ sur $[0;2\pi]$. \hfill (1,5 pt)
  \item Determiner les extremums de $g$ sur $\mathbb{R}$. \hfill (1 pt)
  \item Resoudre $g(x) = -1$ sur $[0;2\pi]$. \hfill (0,5 pt)
\end{enumerate}
